% ============================================================
% M3 S2 导学件·波1 补位1glm 片件 —— 课时03（2.2.1 直线的方向向量与法向量·下）
% 生成：M3 成卷轮 S2 补位1glm 接盘续作｜2026-09-14｜编译 cwd＝本目录（中文路径禁 \input 绝对路径）
%   xelatex main-true.tex ×2 → 含详解印本；xelatex main-false.tex ×2 → 纯题版（单源双本）
% 输入链（全只读）：题面库 课时03-方向向量与法向量.md（题面逐字装配）
%   ＋ 课时03-方向向量与法向量-答案侧.md（值逐字回嵌）＋ 定稿/批A-课时03-方向向量与法向量.md
%   （详解回嵌源 §一/§二，含 0913命制入槽 E12~E16【亲算落账】）；toolchain＝qp-m3.sty
%   （md5 7c3930362be8a0a2bdf21bbf8ac16573，本地挂载副本，禁改；破损即停）。
% 母版体例：承《课时01-坐标法/母版体例说明.md》——课前预习（知识点填空＋判断正误＋书写位，
%   零键零答案）→ 课中探究（3 考点×例＋变式＝练槽 16 键；本课时拓展槽空置·批A§三 在案）
%   → 课堂检测 G5（3单选＋2填空）→ 尾块。
% 键账：件内 ansblock 键集 ≡ 题面库片键集 ≡ manifest 键序（21 键＝E16＋G5，零缺漏零多余）；
%   印面号 1—16＝考点探究（例/变式装配序）＋17—21＝课堂检测 G1—G5；键↔号映射见
%   值台账-课时03.json。多选 0 道（本片题型配比：单选2＋填空5＋解答14）。
% 括线判据（承重墙禁放宽）：块内容估高＞8 行（wlen/23 栏宽字口径，逐键读数入台账）→ 括线模。
% 门谱读数：见 过程件 工作区/_tmpM3S2补位1glm_0914/（编译 log、门谱输出、目验）。
% ============================================================
\documentclass[fontset=none]{ctexart}
\usepackage{qp-m3}
% —— 印面逐字符号路由补钉（探针实证 0914：书宋缺 U+2212，▱ 诸字库皆缺→NSC 子块；
%     禁改 sty 本体保 md5；本片 ₀₁θ∈∞√⊥ 等探针全渲染，⟺ 缺字走 \iff 换算） ——
\xeCJKDeclareCharClass{Default}{"2212}%
\setCJKfamilyfont{nscblk}[Path=C:/提示词/工作区/字替对照-0909/variantF/fonts/,BoldFont=NSC-w600.ttf]{NSC-w500.ttf}%
\xeCJKDeclareSubCJKBlock{nscblk}{"25B1}%
\newcommand{\pxparallelogram}{{\CJKfamily{nscblk}▱}}%
% —— 断行微调（纯题档/详解档三零门：CJK 胶收缩 0.01→0.05em；伸展分量本件按数学密集
%     详解承课时02 实测值 0.25em——胶点稀的公式行需更大伸展量才能撑满栏宽，
%     否则 Underfull \hbox；只动伸缩分量不动自然宽度，版面纹理零漂） ——
\xeCJKsetup{CJKglue={\hskip 0.10em plus 0.25em minus 0.05em}}%
% —— 页脚件名（qp-headfoot 复用入口） ——
\renewcommand{\qpzhangming}{第二章\quad 平面解析几何}%
\renewcommand{\qpjianming}{导学件}%

\begin{document}

\zhangtitle{第二章\quad 平面解析几何}
\jietitle{2.2.1 直线的方向向量与法向量（第2课时）}
\mubiaoline
\mubiaomu{1}{理解直线的方向向量的概念，掌握求方向向量与单位方向向量的方法，会由方向向量求直线的斜率；}
\mubiaomu{2}{理解直线的法向量的概念，掌握方向向量与法向量的互化，会用点法式表示直线的方程；}
\mubiaomu{3}{综合运用方向向量与法向量讨论两直线的位置关系，体会向量方法在解析几何中的工具价值.}

\vspace{1.7mm}
\begin{multicols}{2}
\emergencystretch=1em

% ============================================================
% 课前预习（知识点填空＋判断正误＝件侧自教材 §2.2.1 下半节组织，非题库题；零键零答案，
% \kongbai 空线两档等形不泄答）
% ============================================================
\huaxing{课}{前}{预}{习}{知识导学\quad 素养初识}

\zsd{一}{直线的方向向量}

\tiaomu{1}{直线 \(l\) 上的向量\kongbai{}称为直线 \(l\) 的方向向量；一条直线的方向向量有\kongbai{}个，与直线 \(l\) 平行的任意非零向量都是它的方向向量．\zhuzhu{方向向量只论"共线"不论"起点"，非零是底线；给定一条直线，方向向量不唯一，彼此相差一个非零倍数.}}

\tiaomu{2}{若直线 \(l\) 的方向向量为 \((u,v)\) 且 \(u\neq 0\)，则直线 \(l\) 的斜率 \(k=\)\kongbai{}；若 \(u=0\)，则直线 \(l\) \kongbai{}．\zhuzhu{方向向量与斜率互化是本课时的主工具：横坐标作分母，先查零再作除.}}

\tiaomu{3}{当 \((u,v)\) 满足 \(u^{2}+v^{2}=1\) 时，称其为直线的\kongbai{}；直线 \(l\) 的单位方向向量有\kongbai{}个．\zhuzhu{单位方向向量＝方向向量除以自身长度，"单位"锁定长度为 1，方向相同或相反各一个.}}

\zsd{二}{直线的法向量}

\tiaomu{1}{与直线 \(l\) \kongbai{}的非零向量称为直线 \(l\) 的法向量；若直线 \(l\) 的一个方向向量为 \((u,v)\)，则 \((-v,u)\) 是 \(l\) 的一个\kongbai{}．\zhuzhu{方向向量绕任一点旋转 \(90^{\circ}\) 即得法向量，数量积为零是验算法向量的标准动作.}}

\tiaomu{2}{设直线 \(l\) 过点 \(P(x_0,y_0)\)，其一个法向量为 \((u,v)\)，点 \(Q(x,y)\) 在 \(l\) 上当且仅当 \kongbai{}，这就是直线的点法式方程．\zhuzhu{点法式 \(u(x-x_0)+v(y-y_0)=0\) 与斜率不存在（\(v=0\)）的竖直线兼容，是含竖直线情形的统一写法.}}

\zhentib{(1)}{一条直线有且只有一个方向向量．}

\zhentib{(2)}{零向量也可以作为直线的法向量．}

\zhentib{(3)}{若 \((2,1)\) 是直线 \(l\) 的一个方向向量，则 \((1,-2)\) 是 \(l\) 的一个法向量．}

\zhentib{(4)}{倾斜角为 \(90^{\circ}\) 的直线没有方向向量．}

% ============================================================
% 课中探究（考点探究 3 考点×例＋变式＝练槽 16 键，印面号1—16；题后紧跟答案＝ansblock
% 内嵌，值照答案侧逐字，详解承批A §二 回嵌＋印面清洗：provenance 括注剔除，
% 教学性括注（辨析/注意/合理性）保留；纯文本数学式→\(...\) 换算入台账清洗注记）
% ============================================================
\huaxing{课}{中}{探}{究}{考点探究\quad 素养小结}

% ---------- 探究点一 方向向量的概念与运算 ----------
\tjdnr{一}{方向向量的概念与运算}{例\textbf{1}}{简单(知识点一)}{已知直线\(l\)经过点\(A(-2,0)\)与\(B(-5,3)\)，求直线\(l\)的一个方向向量、斜率\(k\)与倾斜角\(\theta\)．}
\xiexwei{10mm}

\begin{ansblock}[2章-练-课时03-E3]
% ans:2章-练-课时03-E3
\ansitem{1}{方向向量如(1,−1)（答案不唯一）；k＝−1；θ＝135°}
\ansnote{详解}{\(\overrightarrow{AB}=(-5+2,\ 3-0)=(-3,3)\)是直线\(l\)的一个方向向量（与之共线的非零向量如\((1,-1)\)也可）．斜率\(k=(3-0)/(-5-(-2))=3/(-3)=-1\)（也可由方向向量\((1,-1)\)得\(k=-1/1=-1\)）．由\(\tan\theta=-1\)且\(\theta\in[0^{\circ},180^{\circ})\)，得\(\theta=135^{\circ}\)．}
\end{ansblock}

\liB{变式\textbf{2}}{简单(知识点一)}{}{已知直线\(l\)经过点\(P(3,m)\)和点\(Q(m,-2)\)，直线\(l\)的方向向量为\((2,4)\)，则直线\(l\)的斜率为\kongbai{}，实数\(m\)的值为\kongbai{}．}

\begin{ansblock}[2章-练-课时03-E1]
% ans:2章-练-课时03-E1
\ansitem{2}{2；4/3}
\ansnote{详解}{由方向向量\((2,4)\)得直线\(l\)的斜率\(k=4/2=2\)．由两点斜率公式\(k=(m-(-2))/(3-m)=(m+2)/(3-m)\)（\(m\neq 3\)），令\((m+2)/(3-m)=2\)，得\(m+2=6-2m\)，\(3m=4\)，\(m=4/3\)．故斜率为\(2\)，\(m=4/3\)．}
\end{ansblock}

\liB{变式\textbf{3}}{简单(知识点一)}{}{（1）设坐标平面内三点\(A(m,-m-3)\)，\(B(2,m-1)\)，\(C(-1,4)\)，若直线\(AC\)的斜率等于直线\(BC\)的斜率的\(3\)倍，求实数\(m\)的值；（2）已知直线\(l_1\)的方向向量为\(\overrightarrow{n}=(2,1)\)，直线\(l_2\)的倾斜角是直线\(l_1\)倾斜角的\(2\)倍，求直线\(l_2\)的斜率．}
\xiexwei{12mm}

{\ansblockgrayfalse
\begin{ansblock}[2章-练-课时03-E2]
% ans:2章-练-课时03-E2
\ansitem{3}{(1) 1或2；(2) 4/3}
\ansnote{详解}{(1) \(k_{AC}=[4-(-m-3)]/(-1-m)=(m+7)/(-m-1)\)（\(m\neq -1\)），\(k_{BC}=(4-m+1)/(-1-2)=(5-m)/(-3)\)．由\(k_{AC}=3k_{BC}\)：\((m+7)/(-m-1)=3\cdot(5-m)/(-3)=m-5\)，去分母得\(m+7=(m-5)(-m-1)=-m^{2}+4m+5\)，即\(m^{2}-3m+2=0\)，解得\(m=1\)或\(m=2\)，经检验均使两斜率存在，符合题意．\par\noindent (2) 设\(l_1\)的倾斜角为\(\alpha\)，则\(l_2\)的倾斜角为\(2\alpha\)．由方向向量\(\overrightarrow{n}=(2,1)\)得\(\tan\alpha=1/2\)，故\(l_2\)的斜率\(\tan 2\alpha=2\tan\alpha/(1-\tan^{2}\alpha)=1/(1-1/4)=4/3\)．}
\end{ansblock}}

\liB{变式\textbf{4}}{简单(知识点一)}{}{已知\(A(-1,-3)\)，\(B(0,-1)\)，\(C(1,2)\)，\(D(3,5)\)，则\(A\)，\(B\)，\(C\)共线吗？\(A\)，\(B\)，\(D\)呢？}
\xiexwei{10mm}

\begin{ansblock}[2章-练-课时03-E4]
% ans:2章-练-课时03-E4
\ansitem{4}{A、B、C不共线；A、B、D共线．}
\ansnote{详解}{\(\overrightarrow{AB}=(0+1,\ -1+3)=(1,2)\)，\(\overrightarrow{AC}=(1+1,\ 2+3)=(2,5)\)．因为\(1\times 5-2\times 2=1\neq 0\)，\(\overrightarrow{AB}\)与\(\overrightarrow{AC}\)不共线，且两向量有公共起点\(A\)，所以\(A\)、\(B\)、\(C\)不共线．\(\overrightarrow{AD}=(3+1,\ 5+3)=(4,8)=4\overrightarrow{AB}\)，\(\overrightarrow{AB}\)与\(\overrightarrow{AD}\)共线且有公共起点\(A\)，所以\(A\)、\(B\)、\(D\)共线．}
\end{ansblock}

\liB{变式\textbf{5}}{简单(知识点一)}{}{判断命题"如果\(A\)，\(B\)，\(C\)是平面直角坐标系中的三个不同的点，则这三点共线的充要条件是\(\overrightarrow{AB}\)与\(\overrightarrow{BC}\)共线"的真假．}
\xiexwei{10mm}

\begin{ansblock}[2章-练-课时03-E5]
% ans:2章-练-课时03-E5
\ansitem{5}{真命题．证明见详解．}
\ansnote{详解}{必要性：若\(A\)、\(B\)、\(C\)共线，则直线上的向量\(\overrightarrow{AB}\)、\(\overrightarrow{BC}\)及与直线平行的向量都是该直线的方向向量，故\(\overrightarrow{AB}\)与\(\overrightarrow{BC}\)共线．\par\noindent 充分性：若\(\overrightarrow{AB}\)与\(\overrightarrow{BC}\)共线，则直线\(AB\)与直线\(BC\)的方向相同或相反；又两直线过公共点\(B\)，由过一点且方向确定的直线唯一，直线\(AB\)与直线\(BC\)是同一条直线，故\(A\)、\(B\)、\(C\)三点共线．综上命题为真．}
\end{ansblock}

\liB{变式\textbf{6}}{简单(知识点一)}{}{直线\(l\)经过\(A(1,2)\)、\(B(3,6)\)两点，则\(l\)的一个方向向量可以是\nobreak(\kongwei)}

\bindopt {\fontsize{10.5pt}{\qpTJgLeadLiti}\selectfont \makebox[\dimexpr0.5\linewidth-1em\relax][l]{A．\((2,1)\)}\makebox[\dimexpr0.5\linewidth-1em\relax][l]{B．\((1,2)\)}\par}
\bindopt {\fontsize{10.5pt}{\qpTJgLeadLiti}\selectfont \makebox[\dimexpr0.5\linewidth-1em\relax][l]{C．\((-2,1)\)}\makebox[\dimexpr0.5\linewidth-1em\relax][l]{D．\((1,-2)\)}\par}

\begin{ansblock}[2章-练-课时03-E12]
% ans:2章-练-课时03-E12
\ansitem{6}{B}
\ansnote{详解}{\(\overrightarrow{AB}=(3-1,\ 6-2)=(2,4)=2(1,2)\)，故\((1,2)\)为\(l\)的一个方向向量，选B．（辨析：A 中\((2,4)=k(2,1)\)无解，与\(\overrightarrow{AB}\)不成比例；C 中\((-2,1)\cdot(2,4)=-4+4=0\)，是\(l\)的法向量方向而非方向向量；D 中\((1,-2)\)与\(\overrightarrow{AB}\)不成比例，均排除．）}
\end{ansblock}

\liB{变式\textbf{7}}{中档(知识点一)}{}{已知直线\(l\)：\((m+1)x+2y-6=0\)的一个方向向量为\(v=(2,\ m-1)\)，求实数\(m\)的值．}
\xiexwei{10mm}

\begin{ansblock}[2章-练-课时03-E15]
% ans:2章-练-课时03-E15
\ansitem{7}{m=0}
\ansnote{详解}{直线\(Ax+By+C=0\)的方向向量可取\((B,\ -A)\)，故\(l\)的方向向量\(t=(2,\ -(m+1))\)．\(t\parallel v\)（横坐标已同为\(2\)）\(\Rightarrow -(m+1)=m-1\)\(\Rightarrow 2m=0\)\(\Rightarrow m=0\)．（验证：\(m=0\)时\(l\)为\(x+2y-6=0\)，方向向量\((2,-1)=v\)；又法向量\((1,2)\)，\(v\cdot n=2-2=0\)，两法一致．）}
\end{ansblock}

\xiaojie{①识别：方向向量题先定"非零＋共线"两要素，横坐标作分母求斜率先查零.②易错：单位方向向量有"同向、反向"两个，只写一个易漏；三点共线须两向量共线且有公共点.③通法：判共线用坐标行列式\(x_1y_2-x_2y_1=0\)，求参代入后回检分母与题设．}

% ---------- 探究点二 法向量的概念与应用 ----------
\tjdnr{二}{法向量的概念与应用}{例\textbf{8}}{简单(知识点二)}{向量\((x_0,y_0)\)与\((-y_0,x_0)\)中，如果其中一个为直线\(l\)的一个方向向量，则另一个一定是直线\(l\)的一个法向量吗？}
\xiexwei{10mm}

{\ansblockgrayfalse
\begin{ansblock}[2章-练-课时03-E8]
% ans:2章-练-课时03-E8
\ansitem{8}{一定．理由见详解．}
\ansnote{详解}{设\((x_0,y_0)\)是\(l\)的一个方向向量．由方向向量非零知\((x_0,y_0)\neq(0,0)\)，从而\((-y_0,x_0)\neq(0,0)\)．又\((x_0,y_0)\cdot(-y_0,x_0)=-x_0y_0+y_0x_0=0\)，即\((-y_0,x_0)\)与\(l\)的方向向量垂直，于是以\((-y_0,x_0)\)为方向向量的有向线段所在直线与\(l\)垂直，由法向量的定义，\((-y_0,x_0)\)是\(l\)的一个法向量．反之，若\((-y_0,x_0)\)为\(l\)的方向向量，同理\((-y_0,x_0)\neq(0,0)\)即\((x_0,y_0)\neq(0,0)\)，且两向量数量积为\(0\)，故\((x_0,y_0)\)一定是\(l\)的法向量．所以结论成立．}
\end{ansblock}}

\liB{变式\textbf{9}}{中档(知识点二)}{}{已知三点\(A(1,4)\)，\(B(3,3)\)，\(C(4,5)\)．求证：\(AB\perp BC\)（要求用直线的方向向量与法向量的语言完成）．}
\xiexwei{10mm}

\begin{ansblock}[2章-练-课时03-E11]
% ans:2章-练-课时03-E11
\ansitem{9}{证明见详解．}
\ansnote{详解}{\(\overrightarrow{AB}=(3-1,\ 3-4)=(2,-1)\)，故直线\(AB\)的一个方向向量为\(a=(2,-1)\)；\(\overrightarrow{BC}=(4-3,\ 5-3)=(1,2)\)．取\(v=(1,2)\)，验证\(v\)与\(a\)垂直：\(v\cdot a=1\times 2+2\times(-1)=0\)，且\(v\neq 0\)，所以\(v\)是直线\(AB\)的一个法向量．而\(\overrightarrow{BC}=v\)，即直线\(BC\)的方向向量就是直线\(AB\)的法向量，由法向量的定义，直线\(BC\)与直线\(AB\)垂直，故\(AB\perp BC\)．}
\end{ansblock}

\liB{变式\textbf{10}}{中档(知识点二)}{}{已知直线\(l_1\)：\(3x-y+1=0\)与\(l_2\)：\(ax+2y-6=0\)．(1) 写出\(l_1\)的一个法向量；(2) 若\(l_1\perp l_2\)，求\(a\)的值，并用法向量语言说明理由．}
\xiexwei{10mm}

\begin{ansblock}[2章-练-课时03-E14]
% ans:2章-练-课时03-E14
\ansitem{10}{(1) n₁=(3,−1)（任一非零倍数均可）　(2) a=2/3}
\ansnote{详解}{(1) 由一般式系数直接读出\(n_1=(3,-1)\)．(2) \(n_2=(a,2)\)．同一平面内每条直线的法向量与其方向向量互相垂直，将两直线绕任一点同旋\(90^{\circ}\)，方向向量恰转到各自法向量方向且夹角不变，故\(l_1\perp l_2\)\(\iff n_1\perp n_2\)\(\iff n_1\cdot n_2=0\)\(\iff 3a+(-1)\times 2=0\)\(\iff a=2/3\)．}
\end{ansblock}

\liB{变式\textbf{11}}{简单(知识点二)}{}{已知直线\(l\)的一个法向量为\(n=(2,-1)\)，且\(l\)过点\(P(3,1)\)，则\(l\)的方程为\kongbai{}．}

\begin{ansblock}[2章-练-课时03-E13]
% ans:2章-练-课时03-E13
\ansitem{11}{2x−y−5=0}
\ansnote{详解}{法向量\((2,-1)\)\(\Rightarrow l\)具形式\(2x-y+C=0\)．代\(P(3,1)\)：\(6-1+C=0\)\(\Rightarrow C=-5\)．故\(2x-y-5=0\)．（验证：\(P\)代入\(6-1-5=0\)；\(l\)的方向向量\((1,2)\)，\(n\cdot(1,2)=2-2=0\)．）}
\end{ansblock}

\xiaojie{①识别：法向量两要素"非零＋与直线垂直"，与方向向量相差一个\(90^{\circ}\)旋转.②易错：法向量不唯一，"读出\((A,B)\)即可"与"写出全部"是两回事；证垂直先译成"方向向量×法向量"语言.③通法：点法式\(u(x-x_0)+v(y-y_0)=0\)代点定\(C\)，回代验零收尾．}

% ---------- 探究点三 方向向量、法向量与直线方程 ----------
\tjdnr{三}{方向向量、法向量与直线方程}{例\textbf{12}}{中档(知识点三)}{（1）如果直线\(l\)过点\(P(-1,-2)\)，且直线\(l\)的方向向量为\(a=(1,-2)\)，求直线\(l\)的方程；（2）如果直线\(l\)过点\(P(x_0,y_0)\)，且直线\(l\)的方向向量为\(a=(u,v)\)，求直线\(l\)的方程．}
\xiexwei{12mm}

{\ansblockgrayfalse
\begin{ansblock}[2章-练-课时03-E6]
% ans:2章-练-课时03-E6
\ansitem{12}{(1) y＝−2x−4（即2x＋y＋4＝0）；(2) 当u≠0时为y−y₀＝(v/u)(x−x₀)；当u＝0时为x＝x₀．}
\ansnote{详解}{(1) 方向向量\(a=(1,-2)\)的横坐标\(1\neq 0\)，故斜率\(k=-2/1=-2\)，由点斜式得\(y-(-2)=-2[x-(-1)]\)，即\(y=-2x-4\)．\par\noindent (2) 设\(Q(x,y)\)是直线\(l\)上任意一点，则\(\overrightarrow{PQ}=(x-x_0,\ y-y_0)\)与方向向量\(a=(u,v)\)共线，得\(v(x-x_0)-u(y-y_0)=0\)．当\(u\neq 0\)时化为\(y-y_0=(v/u)(x-x_0)\)；当\(u=0\)时（此时\(v\neq 0\)）得\(x-x_0=0\)，即\(x=x_0\)（斜率不存在的竖直线）．}
\end{ansblock}}

\liB{变式\textbf{13}}{中档(知识点三)}{}{（1）如果直线\(l\)过点\(P(1,3)\)，且直线\(l\)的法向量为\(a=(-3,1)\)，求直线\(l\)的方程；（2）如果直线\(l\)过点\(P(x_0,y_0)\)，且直线\(l\)的法向量为\(a=(u,v)\)，求直线\(l\)的方程．}
\xiexwei{12mm}

{\ansblockgrayfalse
\begin{ansblock}[2章-练-课时03-E7]
% ans:2章-练-课时03-E7
\ansitem{13}{(1) y＝3x；(2) u(x−x₀)＋v(y−y₀)＝0，即ux＋vy−ux₀−vy₀＝0．}
\ansnote{详解}{(1) 设\(Q(x,y)\)是直线\(l\)上任意一点，则\(\overrightarrow{PQ}=(x-1,\ y-3)\)与法向量\(a=(-3,1)\)垂直，\(a\cdot\overrightarrow{PQ}=0\)，即\(-3(x-1)+1\cdot(y-3)=0\)，化简得\(y=3x\)．\par\noindent (2) 同法：\(a\cdot\overrightarrow{PQ}=u(x-x_0)+v(y-y_0)=0\)（\(u\)、\(v\)不全为\(0\)，因法向量是非零向量）．这就是直线的点法式方程；当\(v\neq 0\)时可化为斜截形\(y=-(u/v)x+(ux_0+vy_0)/v\)，斜率\(k=-u/v\)（与"方向向量\((v,-u)\)读斜率\(-u/v\)"一致）．}
\end{ansblock}}

\liB{变式\textbf{14}}{中档(知识点三)}{}{已知\(P\)是直线\(l\)上一点，且\(n\)是直线\(l\)的一个方向向量，分别根据下列条件求直线\(l\)的方程：（1）\(P(3,-5)\)，\(n=(1,2)\)；（2）\(P(0,5)\)，\(n=(3,-4)\)．}
\xiexwei{10mm}

\begin{ansblock}[2章-练-课时03-E9]
% ans:2章-练-课时03-E9
\ansitem{14}{(1) y＝2x−11；(2) y＝−(4/3)x＋5（即4x＋3y−15＝0）}
\ansnote{详解}{(1) \(n=(1,2)\)横坐标不为\(0\)，\(k=2/1=2\)，由点斜式\(y+5=2(x-3)\)，即\(y=2x-11\)．\par\noindent (2) \(n=(3,-4)\)横坐标不为\(0\)，\(k=-4/3\)，由点斜式\(y-5=-(4/3)(x-0)\)，即\(y=-(4/3)x+5\)．}
\end{ansblock}

\liB{变式\textbf{15}}{中档(知识点三)}{}{已知\(P\)是直线\(l\)上一点，且\(v\)是直线\(l\)的一个法向量，分别根据下列条件求直线\(l\)的方程：（1）\(P(1,2)\)，\(v=(3,-4)\)；（2）\(P(-1,2)\)，\(v=(3,4)\)．}
\xiexwei{10mm}

\begin{ansblock}[2章-练-课时03-E10]
% ans:2章-练-课时03-E10
\ansitem{15}{(1) 3x−4y＋5＝0；(2) 3x＋4y−5＝0}
\ansnote{详解}{(1) 设\(Q(x,y)\)在\(l\)上，\(v\cdot\overrightarrow{PQ}=3(x-1)-4(y-2)=0\)，化简得\(3x-4y+5=0\)．\par\noindent (2) 同法：\(3(x+1)+4(y-2)=0\)，化简得\(3x+4y-5=0\)．}
\end{ansblock}

\liB{变式\textbf{16}}{简单(知识点三)}{}{直线\(l\)过点\(P(0,1)\)，且一个方向向量为\(v=(1,\sqrt{3})\)，则\(l\)在\(x\)轴上的截距为\kongbai{}．}

\begin{ansblock}[2章-练-课时03-E16]
% ans:2章-练-课时03-E16
\ansitem{16}{−√3/3}
\ansnote{详解}{斜率\(k=\sqrt{3}/1=\sqrt{3}\)，\(l\)：\(y=\sqrt{3}x+1\)．令\(y=0\)：\(x=-1/\sqrt{3}=-\sqrt{3}/3\)．（验证：点\((-\sqrt{3}/3,\ 0)\)代入\(\sqrt{3}\times(-\sqrt{3}/3)+1=-1+1=0\)，斜率与\(v\)一致．）}
\end{ansblock}

\xiaojie{①识别："向量定方程"两条路——方向向量读斜率走点斜式，法向量走点法式，竖直线（横坐标为零/\(v=0\)）点法式兜底.②易错：由\(Ax+By+C=0\)读方向向量取\((B,-A)\)、法向量取\((A,B)\)，两套系数勿混.③通法：求出方程后代已知点回验，向量与直线平行（垂直）用数量积为零复核．}

% ============================================================
% 课堂检测（G5＝导学槽 5 键，印面号 17—21；3单选＋2填空，值照答案侧逐字）
% ============================================================
\begin{ketang}{本组共5题（第17—21题），17—19为单选，20—21为填空．}

\jiancestem{{\fontsize{11.4pt}{14pt}\selectfont\heihao {\numboldjian 17}．}\hspace{4.1mm}\tieside{简单}\hspace{2.2mm}已知直线\(l\)经过点\(P(1,2)\)和点\(Q(-2,-2)\)，则直线\(l\)的单位方向向量为\nobreak(\kongwei)}

\bindopt {\fontsize{10.5pt}{\qpTJgLeadPingjia}\selectfont \makebox[\dimexpr0.5\linewidth-1em\relax][l]{A．\((-3,-4)\)}\makebox[\dimexpr0.5\linewidth-1em\relax][l]{B．\((-3/5,-4/5)\)}\par}
\bindopt {\fontsize{10.5pt}{\qpTJgLeadPingjia}\selectfont \makebox[\dimexpr0.5\linewidth-1em\relax][l]{C．\((\pm 3/5,\pm 4/5)\)}\makebox[\dimexpr0.5\linewidth-1em\relax][l]{D．\(\pm(3/5,4/5)\)}\par}

\begin{ansblock}[2章-导-课时03-G1]
% ans:2章-导-课时03-G1
\ansitem{17}{D}
\ansnote{详解}{直线\(l\)的一个方向向量为\(\overrightarrow{PQ}=(-2-1,\ -2-2)=(-3,-4)\)，\(|\overrightarrow{PQ}|=\sqrt{(-3)^{2}+(-4)^{2}}=5\)，故单位方向向量为\(\pm\overrightarrow{PQ}/|\overrightarrow{PQ}|=\pm(1/5)(-3,-4)=\pm(3/5,4/5)\)（方向相同或相反的两个单位向量都是方向向量）．C 写成"\(\pm\)"作用于每个分量的四种组合，其中\((3/5,-4/5)\)不与\(\overrightarrow{PQ}\)共线，故C错．选D．}
\end{ansblock}

\jiancestem{{\fontsize{11.4pt}{14pt}\selectfont\heihao {\numboldjian 18}．}\hspace{4.1mm}\tieside{简单}\hspace{2.2mm}若过点\(P(3,2m)\)和点\(Q(-m,2)\)的直线与方向向量为\(a=(-5,5)\)的直线平行，则实数\(m\)的值是\nobreak(\kongwei)}

\bindopt {\fontsize{10.5pt}{\qpTJgLeadPingjia}\selectfont \makebox[\dimexpr0.5\linewidth-1em\relax][l]{A．\(1/3\)}\makebox[\dimexpr0.5\linewidth-1em\relax][l]{B．\(-1/3\)}\par}
\bindopt {\fontsize{10.5pt}{\qpTJgLeadPingjia}\selectfont \makebox[\dimexpr0.5\linewidth-1em\relax][l]{C．\(2\)}\makebox[\dimexpr0.5\linewidth-1em\relax][l]{D．\(-2\)}\par}

\begin{ansblock}[2章-导-课时03-G2]
% ans:2章-导-课时03-G2
\ansitem{18}{B}
\ansnote{详解}{\(\overrightarrow{PQ}=(-m-3,\ 2-2m)\)．两直线平行\(\iff\)方向向量共线（此处即向量共线，注意\(P\)、\(Q\)不重合需\(2m\neq 2\)即\(m\neq 1\)），由共线条件\((-m-3)\times 5-(2-2m)\times(-5)=0\)，得\(-5m-15+10-10m=0\)，\(-15m-5=0\)，\(m=-1/3\)．经检验\(m=-1/3\)符合题意．选B．}
\end{ansblock}

\jiancestem{{\fontsize{11.4pt}{14pt}\selectfont\heihao {\numboldjian 19}．}\hspace{4.1mm}\tieside{简单}\hspace{2.2mm}已知直线\(l\)的一个法向量为\(v=(2,-4)\)，则直线\(l\)的斜率为\nobreak(\kongwei)}

\bindopt {\fontsize{10.5pt}{\qpTJgLeadPingjia}\selectfont \makebox[\dimexpr0.5\linewidth-1em\relax][l]{A．\(1/2\)}\makebox[\dimexpr0.5\linewidth-1em\relax][l]{B．\(-1/2\)}\par}
\bindopt {\fontsize{10.5pt}{\qpTJgLeadPingjia}\selectfont \makebox[\dimexpr0.5\linewidth-1em\relax][l]{C．\(2\)}\makebox[\dimexpr0.5\linewidth-1em\relax][l]{D．\(-2\)}\par}

\begin{ansblock}[2章-导-课时03-G3]
% ans:2章-导-课时03-G3
\ansitem{19}{A}
\ansnote{详解}{法向量\(v=(2,-4)\)垂直于直线\(l\)，故\(l\)的一个方向向量为与\(v\)垂直的非零向量\(a=(4,2)\)（因\(v\cdot a=2\times 4+(-4)\times 2=0\)），其横坐标不为\(0\)，斜率\(k=2/4=1/2\)．故选A．（另解：法向量为\((A,B)\)的直线可写成\(Ax+By+C=0\)形，\(2x-4y+C=0\)即\(y=x/2+C/4\)，\(k=1/2\)．）}
\end{ansblock}

\jiancestem{{\fontsize{11.4pt}{14pt}\selectfont\heihao {\numboldjian 20}．}\hspace{4.1mm}\tieside{简单}\hspace{2.2mm}已知点\(A(0,-4)\)，\(B(3,0)\)，则直线\(AB\)的一个方向向量为\kongbai{}，线段\(AB\)的长度为\kongbai{}．}
\xiexwei{8mm}

\begin{ansblock}[2章-导-课时03-G4]
% ans:2章-导-课时03-G4
\ansitem{20}{(3,4)（答案不唯一）；5}
\ansnote{详解}{直线\(AB\)的一个方向向量为\(\overrightarrow{AB}=(3-0,\ 0+4)=(3,4)\)（与之平行的任何非零向量均可）；线段\(AB\)的长度为\(|AB|=\sqrt{3^{2}+4^{2}}=5\)．}
\end{ansblock}

\jiancestem{{\fontsize{11.4pt}{14pt}\selectfont\heihao {\numboldjian 21}．}\hspace{4.1mm}\tieside{中档}\hspace{2.2mm}已知经过坐标平面内\(A(1,2)\)，\(B(-2,2m-1)\)两点的直线的方向向量为\((1,\sin\alpha)\)，则实数\(m\)的取值范围为\kongbai{}．}
\xiexwei{8mm}

\begin{ansblock}[2章-导-课时03-G5]
% ans:2章-导-课时03-G5
\ansitem{21}{[0,3]}
\ansnote{详解}{方向向量的横坐标为\(1\neq 0\)，故直线\(AB\)的斜率存在，设为\(k\)．由方向向量\((1,\sin\alpha)\)知\(k=\sin\alpha/1=\sin\alpha\in[-1,1]\)；又由两点坐标\(k=(2m-1-2)/(-2-1)=(3-2m)/3\)．所以\(-1\leqslant(3-2m)/3\leqslant 1\)，解得\(0\leqslant m\leqslant 3\)．故\(m\)的取值范围为\([0,3]\)．}
\end{ansblock}

\end{ketang}

% ---- 尾块（\tailfill 置 \end{multicols} 前；无参＝内容尾随框，件侧不擅自定高） ----
\tailfill

\end{multicols}

\end{document}
